\documentclass{article}
\usepackage{graphicx} 
\usepackage{float}   

\begin{document}

\begin{titlepage}
    \centering
    
    {\Large Universidad Distrital Francisco José de Caldas \par}
    \vspace{1.5cm}
    
    {\Huge \textbf{INFORME: ANÁLISIS DE LAS REACCIONES DE HIDRÓLISIS EN DERIVADOS DE ACILO} \par}
    \vspace{1cm}
    
    {\Large Proyecto Final \par}
    \vspace{3cm}
    
    {\Large Química Orgánica \par}
    \vspace{3cm}
    
    {\large 
    Nombre: Jade Gonzalez, Laura Vega, Paula Martínez, Luna Prieto \par
    Profesor: Oscar Genario Garcia \par
    }
    
    \vfill
    
    {\large Bogotá, Colombia \par}
    {\large Mayo 2026 \par}

\end{titlepage}

\section{Resumen}
El presente documento consolida y analiza las propiedades físicas, químicas y el comportamiento reactivo de diversas familias de compuestos orgánicos, con un enfoque primordial en los procesos hidrolíticos. A través de la revisión sistemática de cinco grupos funcionales ---nitrilos, haluros de acilo, anhídridos de ácido, ésteres y amidas--- se documentan sus reacciones características. El análisis aborda la \textit{sustitución nucleofílica acílica}, considerando las variaciones cinéticas y termodinámicas de los sustratos frente a diferentes medios de reacción (neutros, ácidos y básicos). Asimismo, se integra el mapeo de parámetros fisicoquímicos fundamentales recopilados en el estudio (como el peso molecular, punto de ebullición, densidad, constante de acidez, pKa y toxicidad y LD50). El enfoque central en las reacciones de hidrólisis permite establecer una métrica estándar para evaluar la fortaleza de los grupos salientes y predecir con precisión la estabilidad relativa de los derivados acílicos analizados.


\section{Introducción}
La hidrólisis es una de las reacciones más importantes dentro de la química orgánica, ya que permite la transformación de diferentes compuestos mediante la acción del agua, generando nuevos productos con propiedades químicas distintas. Este proceso ocurre a través de la ruptura de enlaces químicos y puede desarrollarse en medio ácido o básico, dependiendo de las condiciones de reacción del compuesto involucrado, la hidrólisis representa un tema fundamental para comprender el comportamiento y la reactividad de los compuestos orgánicos.                                                                                                                                           En el presente proyecto se estudiaron las reacciones de hidrólisis en distintos grupos funcionales, entre ellos nitrilos, amidas, ésteres, haluros de ácido y anhídridos. A través del análisis de estas reacciones se buscó comprender la manera en que cada grupo funcional responde frente al proceso de hidrólisis, identificando los cambios estructurales que ocurren durante la transformación química y las condiciones necesarias para que dichas reacciones se desarrollen.

El estudio de la hidrólisis resulta fundamental porque este tipo de reacciones tiene una amplia aplicación tanto en el laboratorio como en diferentes procesos industriales. Además, permite comprender la reactividad de los compuestos orgánicos y la formación de sustancias de interés químico como ácidos carboxílicos, alcoholes, amidas y sales orgánicas. Por esta razón, el desarrollo de este proyecto contribuye al fortalecimiento de los conocimientos sobre química orgánica y facilita la relación entre la teoría y las transformaciones químicas observadas en los diferentes grupos funcionales estudiados.

\section{Marco teórico}
El estudio sistemático de las transformaciones moleculares en la química de compuestos oxigenados y nitrogenados exige una comprensión detallada de las interacciones núcleo-electrón. A continuación, se fundamentan las bases conceptuales que rigen los procesos de los derivados de ácido y nitrilos.

Hidrólisis
La hidrólisis es una reacción química fundamental en la que una molécula de agua (H2O) interactúa con un sustrato orgánico, provocando la ruptura de uno o más enlaces covalentes. En este proceso, la molécula de agua se fragmenta en un protón y un anión hidroxilo. Aplicado a los derivados de ácido carboxílico, el proceso hidrolítico implica una sustitución nucleofílica en el grupo acilo, donde el grupo saliente es desplazado por el grupo hidroxilo, revirtiendo el derivado a su forma base: el ácido carboxílico progenitor.

Grupos Funcionales y su Relación con la Hidrólisis.
Los grupos funcionales son determinantes directos de la reactividad. Los compuestos analizados: haluros de acilo, anhídridos, ésteres, amidas y nitrilos, presentan diferencias críticas en la densidad electrónica de su carbono electrofílico (el carbono carbonílico o el carbono del cianuro). La susceptibilidad a la hidrólisis de estos compuestos está intrínsecamente ligada a dos factores: la capacidad de polarización del enlace carbono-grupo saliente y la debilidad del grupo saliente como base conjugada. Cuanto mejor sea el grupo saliente y más deficiente en electrones se encuentre el carbono, más rápida y exotérmica será la hidrólisis.

Hidrólisis de Nitrilos.
Los nitrilos, aunque no son formalmente compuestos carbonílicos, reaccionan análogamente debido al marcado carácter electrofílico del carbono sp polarizado por el nitrógeno. Su hidrólisis transcurre en dos etapas secuenciales: una hidratación inicial que convierte el nitrilo en una amida primaria, y una subsecuente hidrólisis de la amida para formar el ácido carboxílico correspondiente (o su sal).

Hidrólisis de Amidas.
Las amidas representan la familia más estable frente a la hidrólisis. Esto se debe a la significativa estabilización por resonancia proporcionada por el par de electrones no compartidos del nitrógeno, que confiere un carácter de doble enlace al enlace ya que el grupo amiduro es una base sumamente fuerte y, por ende, un pésimo grupo saliente. La hidrólisis de amidas exige condiciones drásticas: reflujo prolongado con ácidos o bases acuosas concentradas.

Hidrólisis de Ésteres. 
Los ésteres ocupan una posición de reactividad intermedia. Su hidrólisis puede ocurrir tanto en medio ácido (proceso reversible que alcanza un equilibrio, regulado por el Principio de Le Châtelier) como en medio básico. La hidrólisis básica, ampliamente conocida como saponificación, es de particular relevancia industrial en la fabricación de jabones.

Hidrólisis de Haluros de Acilo.
Son los derivados más reactivos. La alta electronegatividad del halógeno acentúa la carga parcial positiva sobre el carbono, mientras que, al mismo tiempo, el ion haluro es un excelente grupo saliente debido a su carácter de base extremadamente débil. En consecuencia, los haluros de acilo se hidrolizan de forma casi inmediata y violenta al entrar en contacto incluso con humedad atmosférica o agua fría, formando el ácido carboxílico y el hidrácido.

Hidrólisis de Anhídridos.
Los anhídridos de ácido poseen un grupo saliente moderadamente bueno (un ion carboxilato, el cual está estabilizado por resonancia). Su reactividad es inferior a la de los haluros, pero superior a la de ésteres y amidas. Se hidrolizan gradualmente con agua a temperatura ambiente, o más rápidamente con ligera calefacción, originando dos moléculas del ácido carboxílico correspondiente.

Hidrólisis en Medio Ácido y Medio Básico
Las condiciones del medio modifican el mecanismo de reacción.

Medio Ácido: La reacción se inicia con la protonación del oxígeno del carbonilo (o nitrógeno del cianuro). Esto incrementa drásticamente la vulnerabilidad del carbono central, permitiendo que una molécula neutra de agua y débilmente nucleofílica lo ataque.

Medio Básico: El reactivo principal es el ion hidróxido, un nucleófilo fuerte capaz de atacar directamente al carbonilo sin requerir activación previa. Este medio hace que la reacción de derivados acílicos sea irreversible, ya que el producto final es un anión carboxilato, el cual no es susceptible de ataques nucleofílicos adicionales.

Importancia de la Hidrólisis en el Proyecto (Porqué se analizaron estas reacciones)
El enfoque en las reacciones de hidrólisis dentro del desarrollo de este proyecto analítico no es fortuito; obedece a razones químicas, metodológicas y pedagógicas clave, fuertemente alineadas con la matriz de datos investigada:

Escala Universal de Reactividad: La hidrólisis actúa como el "denominador común" y el patrón de medida perfecto. Al utilizar siempre el mismo reactivo (en medio acuoso), permite comparar experimentalmente la estabilidad de todas las familias funcionales bajo una misma lupa, validando la serie teórica de reactividad, donde los Haluros son los mayores, seguidos de los Anhídridos, luego Ésteres y por últimos Amidas y Nitrilos.
Evaluación Directa de la Fortaleza del Grupo Saliente: Al observar la necesidad de incluir catalizadores y variar las temperaturas (como se detalla en las variables del archivo central), la hidrólisis expone empíricamente la relación inversa entre la basicidad del grupo saliente y la velocidad de reacción.
Justificación de la Matriz de Variables Físico-Químicas: El análisis hidrolítico dota de significado a los parámetros del proyecto. Por ejemplo, entender la hidrólisis explica por qué un compuesto con alta solubilidad y alta constante de acidez puede requerir medios anhidros para su manipulación (caso de los haluros), o por qué ciertas reacciones necesitan tiempos prolongados en solventes fuertemente ionizantes.


\section{Resultados obtenidos}
Durante el análisis de la tabla con enfoque en los datos de reacciones de tipo hidrólisis, se obtuvo una variedad de gráficas con respecto a las propiedades fisicoquímicas como punto de ebullición, peso molecular, densidad, pKa y LD50. Previo a la construcción de estas gráficas, se ejecutó un proceso de depuración y filtrado de datos (remoción de valores nulos o inconsistentes); esto con el fin de garantizar una mayor precisión en el cálculo de los promedios y asegurar la confiabilidad del análisis estadístico visual.
\begin{figure}[H] 
    \centering
    \includegraphics[width=1\textwidth]{cantidad de reacciones de hidrólisis.png}
    \caption{Cantidad de reacciones de hidrólisis por grupo funcional}
    \label{fig:punto_ebullicion}
\end{figure}

Este esquema fue diseñado para observar de manera más estadística la cantidad de datos que se tendrán en cuenta en cada una de las gráficas, siendo evidente donde está la mayor cantidad de estas reacciones dentro de los grupos funcionales que son los nitrilos. 

\begin{figure}[H] 
    \centering
    \includegraphics[width=1\textwidth]{punto de ebullicion productos.png}
    \caption{punto de ebullición del producto}
    \label{fig:punto_ebullicion}
\end{figure}
En esta gráfica podemos observar el punto de ebullición de los productos que pasaron por una reacción de hidrólisis, evaluado en cada grupo funcional. Donde se puede deducir que el mayor punto de ebullición se encuentra dentro de los nitrilos. 

\begin{figure}[H] 
    \centering
    \includegraphics[width=1\textwidth]{punto de ebullición sustratos.png}
    \caption{puntos de ebullición del sustrato}
    \label{fig:punto_ebullicion}
\end{figure}
En esta gráfica del punto de ebullición de los sustratos, nuevamente podemos deducir que el mayor punto de ebullición se encuentra dentro de los nitrilos, lo que podría indicar un equilibrio entre el sustrato y el producto después de la hidrólisis.

Después de mostrar el punto de ebullición tanto del sustrato como del producto en las reacciones de hidrólisis, se pasará a mostrar el peso molecular de las mismas. 

\begin{figure}[H] 
    \centering
    \includegraphics[width=1\textwidth]{peso molecular productos.png}
    \caption{peso molecular de los productos}
    \label{fig:peso molecular}
\end{figure}

Aquí se puede observar que hay cierta estabilidad en los pesos moleculares dentro de algunas de las reacciones de hidrólisis en grupos funcionales como amidas, anhídridos, haluros de acilo y ésteres; caso contrario a los nitrilos, ya que sus productos tienen valores de peso molecular dispersos, lo que podría indicar que sus pesos moleculares no son tan estables con respecto a los otros grupos funcionales.

\begin{figure}[H] 
    \centering
    \includegraphics[width=1\textwidth]{peso molecular sustrato y producto.png}
    \caption{peso molecular de sustratos y productos}
    \label{fig:peso molecular}
\end{figure}
En esta gráfica se puede observar la variabilidad del peso molecular del sustrato y producto dependiendo su grupo funcional; también correlaciones estables que podemos ver en amidas. 

\begin{figure}[H] 
    \centering
    \includegraphics[width=1\textwidth]{densidad.png}
    \caption{Comparación de densidad entre sustrato y producto}
    \label{fig:densidad}
\end{figure}

Dentro de esta gráfica comparativa de análisis en la densidad, se puede observar como hay cierta correlación en los valores de densidad de sustrato y producto en los grupos funcionales de amidas, anhídridos, haluros de acilo y ésteres; mientras que en los nitrilos esto no sucede, ya que sus productos tienden a tener una mayor densidad que sus sustratos, exceptuando el caso excepcional del sustrato que sobrepasa los niveles de densidad de los productos. 


Dentro de los resultados obtenidos se logró sacar únicamente los valores de pKa de los productos, puesto que no habían los suficientes datos de pKa en los sustratos de anhídirdos (esto se explicará de manera breve en la siguiente sección del informe). 

\begin{figure}[H] 
    \centering
    \includegraphics[width=1\textwidth]{nuevo.png}
    \caption{Datos de variabilidad en pKa con respecto a los grupos funcionales.}
    \label{fig:pka}
\end{figure}
En esta sección se puede observar como el sustrato con mayor valor de pKa se encuentra dentro de las amidas, y como los anhídridos permanecen a un nivel muy bajo de pKa con respecto a los demás grupos funcionales; también se observa que en los nitrilos vuelve a haber una dispersión en sus datos. 

Ahora se mostrará finalmente la gráfica de LD50 (toxicidad) 

\begin{figure}[H] 
    \centering
    \includegraphics[width=1\textwidth]{Pkanuevo.png}
    \caption{Toxicidad de sustratos en reacción de hidrólisis de cada grupo funcional.}
    \label{fig:LD50}
\end{figure}

De esta gráfica se puede deducir que la mayor toxicidad está dentro de los sustratos de nitrilos debido debido a la baja docificación que puede llegar a ser letal para un mamífero. Por el contrario, los esteres son los compuestos más seguros.
\section{Análisis de resultados}
A partir de las gráficas realizadas, se recopilaron y contrastaron los datos de reactivos (sustratos) y productos evaluados en función del punto de ebullición, peso molecular y densidad. Este análisis conjunto permitió comprender un modelo de comportamiento predictivo en los derivados de ácidos carboxílicos y su reactividad frente a la hidrólisis.

Principalmente, en el análisis de la Figura \ref{fig:peso molecular} fueron captados los datos de punto de ebullición de los sustratos, donde al contrastarse con los productos de la Figura \ref{fig:peso molecular}, se evidencia que la mayor brecha térmica entre los derivados y sus respectivos productos (sintetizados como ácidos carboxílicos) la presentaban los nitrilos, seguida de los anhídridos, ésteres, haluros de acilo y, por último, las amidas. Estas últimas indicaron una gran similitud en sus temperaturas de ebullición con su ácido homólogo.

Sin embargo, al evaluar la Figura Figura \ref{fig:densidad} de densidades, todos los sustratos presentaron densidades diferentes respecto a las de sus productos, incluyendo a las amidas y sus ácidos carboxílicos, a pesar de su punto de ebullición y su relación con el peso molar. Esto se debe a que las amidas poseen un grupo amino NH2 altamente polar por resonancia con el carbonilo C=O. Al evaluar este par homólogo, sus fuerzas de dispersión de London son equivalentes debido a la similitud en sus masas molares; sin embargo, la amida logra formar una red tridimensional extendida de puentes de hidrógeno. En contraste, el ácido acético, aunque forma puentes de hidrógeno muy fuertes, prefiere unirse en parejas cerradas llamadas dímeros cíclicos (Carey \& Sundberg, 2007). Al confinarse en parejas moleculares independientes, disminuye su cohesión general y, por ende, reduce la necesidad de calor para hervir en comparación con la amida.

Por otro lado, mediante la misma Figura \ref{fig:densidad} se analizó el comportamiento volumétrico de los anhídridos, los cuales destacan como el único grupo funcional capaz de mantener un rango de oscilación definido y geométricamente constante entre el sustrato y el producto. Esto responde directamente a la simetría de su estructura molecular; al reaccionar con el agua, la molécula experimenta una ruptura simétrica que produce exactamente dos moléculas idénticas de ácido carboxílico. Al conservarse la relación de masa y el volumen ocupado en el espacio, la densidad macroscópica no sufre variaciones significativas. De este modo, en la gráfica de densidades el anhídrido es el único que no da un salto y mantiene su rango constante, debido a que en la Figura \ref{fig:peso molecular} de peso molecular, se evidencia una caída simétrica donde el peso de partida equivale a la suma de los dos ácidos carboxílicos generados como producto. Este monitoreo de la masa molecular permite aislar el efecto de las fuerzas de dispersión de London en las tendencias de empaquetamiento, comprobando que las variaciones en las gráficas térmicas y volumétricas se rigen primordialmente por la reconfiguración de los enlaces dirigidos (puentes de hidrógeno) tras la hidrólisis.

En cuanto a los ésteres y nitrilos de partida, estos alcanzan una posición inicial constante en las curvas porque no pueden formar puentes de hidrógeno al carecer de hidrógenos unidos directamente a un oxígeno o a un nitrógeno. Su unión en estado de sustrato depende estrictamente de interacciones dipolo-dipolo (Katritzky et al., 2000); al atraerse con menos fuerza, sus moléculas están más separadas en el espacio, lo que vuelve al líquido menos denso y provoca que se evapore a temperaturas más bajas. Sin embargo, tras el proceso de hidrólisis (Figura \ref{fig:densidad)}), los productos de ambos grupos experimentan un incremento drástico en sus valores de densidad. Este fenómeno se debe a la aparición de los puentes de hidrógeno del ácido carboxílico, los cuales inducen una severa contracción en el volumen libre del líquido, forzando un empaquetamiento mucho más compacto y elevando de manera uniforme la densidad en los productos de ambas familias funcionales debido a la reorganización estructural descrita por los modelos de relación estructura-propiedad.

Una de las propiedades analizadas es el pKa, evaluado en la Figura \ref{fig:pka} a partir de los productos de hidrólisis, ya que los sustratos más reactivos presentan vacíos experimentales (N/A). Los datos muestran valores altos de pKa (y por tanto menor acidez) en los derivados de amidas, seguidos por nitrilos y ésteres. En contraste, los haluros de acilo y anhídridos presentan los menores valores de pKa, asociados a pH bajos y una acidez mucho mayor en solución acuosa (Carey \& Sundberg, 2007). Además, aunque el pKa corresponde al ácido carboxílico final y no depende directamente del derivado de acilo inicial, la alta acidez macroscópica en haluros y anhídridos se intensifica por la liberación de subproductos fuertemente ácidos, como el ácido clorhídrico o ácidos homólogos libres.

Por otra parte, al evaluar la Figura \ref{fig:LD50}, la cual indica los niveles de toxicidad en mamíferos LD50, se observaron las cantidades en miligramos (mg/kg) necesarias para interferir con la salud de un mamífero. Mientras que los ésteres exhiben una dosificación letal más alta, lo cual refleja un perfil de seguridad biológica mucho más alto y seguro, los compuestos altamente reactivos muestran umbrales críticos drásticamente inferiores.

Finalmente, la falta de datos (N/A) en la acidez y toxicidad de los sustratos más reactivos dentro de la plataforma CAS-SciFinder tiene una explicación cinética fundamental. Debido a que los haluros de acilo y los anhídridos poseen grupos salientes muy estables (como el ion cloruro Cl- o el grupo acetato), su energía de activación para la sustitución nucleofílica de acilo es bajísima (Carey \& Sundberg, 2007). Esto desata una hidrólisis instantánea e irreversible al menor contacto con la humedad del ambiente. Esta extrema inestabilidad química impide medir experimentalmente un pKa o una LD50 in vivo de la molécula intacta en soluciones acuosas estándar, ya que el sustrato se destruye antes de poder registrar sus propiedades macroscópicas, obligando al estado del arte a depender de aproximaciones de modelado computacional in silico para estimar el riesgo latente de la sustancia pura.

Finalmente, la falta de datos (N/A) en la acidez y toxicidad de los sustratos más reactivos dentro de la plataforma CAS-SciFinder tiene una explicación cinética fundamental. Debido a que los haluros de acilo y los anhídridos poseen grupos salientes muy estables (como el ión cloruro Cl- o el acetato), su energía de activación para la sustitución nucleofílica de acilo es bajísima (Carey \& Sundberg, 2007). Esto desata una hidrólisis instantánea e irreversible al menor contacto con la humedad del ambiente. Esta extrema inestabilidad química impide medir experimentalmente un pKa o una LD50 in vivo de la molécula intacta en soluciones acuosas estándar, ya que el sustrato se destruye antes de poder registrar sus propiedades macroscópicas, obligando al estado del arte a depender de aproximaciones de modelado computacional in silico para estimar el riesgo latente de la sustancia pura.

\subsection{Comparación y contraste con la teoría}
\subsubsection{Desviación térmica en amidas}

Teóricamente, las amidas presentan altos puntos de ebullición debido a su capacidad de formar puentes de hidrógeno (Galamba, 2015). Sin embargo, en la Figura 1 no ocuparon la posición más alta esperada. Esto puede explicarse porque la mayoría de compuestos analizados corresponden a cadenas alifáticas cortas, donde la baja polarizabilidad electrónica disminuye el efecto de las fuerzas de London y limita el aumento del punto de ebullición.

\subsubsection{Reactividad de haluros de acilo y anhídridos}
La ausencia de datos experimentales estables para haluros de acilo y anhídridos respalda su alta reactividad frente a la hidrólisis. Tanto el ion cloruro como el grupo acetato funcionan como excelentes grupos salientes, reduciendo la energía de activación de la reacción (Bentley, 2011). Esto explica su rápida descomposición incluso en presencia de humedad ambiental y justifica el uso de modelos predictivos in silico para evaluar propiedades como pKa y LD50 (Gleeson, 2008).

\subsubsection{Incremento de densidad luego de hidrolizar los derivados de acilo}
El aumento de densidad observado en ésteres y nitrilos coincide con los modelos de relación estructura-propiedad (QSPR). Después de la hidrólisis, los productos obtenidos pueden formar puentes de hidrógeno más fuertes, generando un empaquetamiento molecular más compacto y un incremento significativo en la densidad del sistema (Katritzky et al., 2000; Katritzky et al., 2003).

\section{Conclusiones}
Existe una relación directa y predecible entre la estructura del grupo funcional del sustrato/producto y su temperatura de ebullición. Las amidas y nitrilos consolidan los puntos de ebullición más altos de toda la muestra debido a la fuerte intensidad de sus fuerzas intermoleculares (redes de puentes de hidrógeno en amidas y fuertes momentos dipolo-dipolo en nitrilos). En contraste, los haluros de acilo y ésteres se ubican en los rangos inferiores debido a interacciones moleculares más débiles.

Las gráficas comparativas bilaterales (sustrato vs. producto) validan visualmente el mecanismo químico de la hidrólisis como una reacción de ruptura. Al analizar los anhídridos, se evidencia de forma estadística una disminución marcada en el peso molecular al pasar de sustrato a producto. Por su parte, los ésteres muestran una estabilización en la distribución de masa de sus productos principales, reflejando que la conversión funcional reorganiza las fuerzas de cohesión sin requerir un colapso masivo en el peso molecular de su cadena principal.
\section{Bibliografías}

Wade, L. G. (2011). Química Orgánica (7ma ed.). Pearson Educación. 
McMurry, J. (2012). Química Orgánica (8va ed.). 
Cengage Learning. Clayden, J., Greeves, N. Warren, S. (2012). Organic Chemistry (2nd ed.). 
Oxford University Press. Bruice, P. Y. (2015). Química Orgánica (7ma ed.). Pearson Educación.

Carey \& Sundberg (2007) — Advanced Organic Chemistry (Libro Guía) 
https://link.springer.com/book/10.1007/978- \- 0-387-44899-3


Katritzky et al. (2000) — QSPR for predicting physical properties. \- https://pubs.acs.org/doi/10.1021/ci9903206


Rumble / CRC Press (2020) — CRC Handbook of Chemistry and Physics
https://www.routledge.com/CRC-Handbook-of- \- Chemistry-and-Physics/book-series/CRCHANCHEPHY

\end{document}